\scanpage{102}\mspage{1}
% archive foliation: 101
Soit $E$ courbe elliptique, sur corps alg clos\inserted{$k$} pour
commencer (de car $\ne2$ comme d'habitude).
On pose, pour $n\in\mathbb N$
\[
(1)\qquad\mathcal F_n=\Gamma(E,n\{0_E\})
\]
par la formule de RR on a
\[
(2)\qquad\operatorname{rg}_k\mathcal F_n=n\qquad\text{si }n\geq1
\]
et
\[
(3)\qquad\mathcal F_0=\mathcal F_1=k.1_E.
\]
On a \deleted{\unsure{une}} sur
\[
(4)\qquad\mathcal F_\infty\overset{\mathrm{def}}=\varinjlim_n\mathcal F_n
\]
une loi de multiplication ass. comm.
\[
(5)\qquad\mathcal F_n.\mathcal F_m\subset\mathcal F_{n+m},
\]
de sorte que $\mathcal F_*=\unsure{\coprod}\mathcal F_n$ devient une
algèbre graduée, dont le $\operatorname{Proj}$ est can. isom.
à $E$. On va déterminer $\mathcal F_*$.
Notons la suite exacte
\[
(6)\qquad
\diagram{\begin{tikzcd}[column sep=small]
0\arrow[r, no head]&\mathcal F_{n-1}\arrow[r,no head]&\mathcal F_n\arrow[r,"\varepsilon_n"]&t_E^{\otimes n}\arrow[r]&0
\end{tikzcd}}
\qquad(n\geq2)
\]
où
\[
(7)\qquad\underline t_E\overset{\mathrm{def}}=T_{E,0_E}.
\]
On a, pour $\delta\in t_E$, une opération
